% Generated from the audited Markdown source; author: Carptopus.
\documentclass[11pt]{article}
\usepackage[a4paper,margin=25mm]{geometry}
\usepackage[T1]{fontenc}
\usepackage[utf8]{inputenc}
\usepackage{lmodern}
\usepackage{microtype}
\usepackage{amsmath,amssymb,amsthm,mathtools}
\usepackage{needspace}
\usepackage{xcolor}
\usepackage[colorlinks=true,linkcolor=blue!55!black,citecolor=blue!55!black,urlcolor=blue!65!black]{hyperref}
\usepackage{fancyhdr}

\newtheorem{theorem}{Theorem}[section]
\newtheorem{lemma}[theorem]{Lemma}
\newtheorem{fact}[theorem]{Fact}
\theoremstyle{remark}
\newtheorem{remark}[theorem]{Remark}
\numberwithin{equation}{section}
\allowdisplaybreaks[2]

\pagestyle{fancy}
\fancyhf{}
\fancyhead[L]{Carptopus}
\fancyhead[R]{Third support regions of binary $RM_2$ codes}
\fancyfoot[C]{\thepage}
\setlength{\headheight}{14pt}
\setlength{\parindent}{0pt}
\setlength{\parskip}{0.42em}
\setlength{\emergencystretch}{4em}
\urlstyle{same}

\title{Exact regions and low-polar-rank recursion\\
in the third support spectrum of binary\\
second-order Reed--Muller codes}
\author{Carptopus\\\href{mailto:carptopus@163.com}{carptopus@163.com}}
\date{30 August 2026}

\hypersetup{
  pdftitle={Exact regions and low-polar-rank recursion in the third support spectrum of binary second-order Reed-Muller codes},
  pdfauthor={Carptopus},
  pdfsubject={Exact regions, finite low-rank recursion, and a forbidden ray in the three-dimensional support spectrum of binary second-order Reed-Muller codes},
  pdfkeywords={Reed-Muller code, higher support spectrum, quadratic Boolean function, Walsh transform, alternating pencil, Fano plane, Kronecker block}
}

\begin{document}
\maketitle

\begin{center}
\fcolorbox{black!35}{black!4}{\parbox{0.88\textwidth}{\centering\small
Public Beta manuscript, version 0.1-beta. The mathematical proof and the
complete integrated manuscript passed project-internal zero-trust audits.
External mathematical review and formal peer review are pending.}}
\end{center}

\begin{abstract}


Let \(\mathcal S_n^{(3)}\) be the set of support sizes of three-dimensional
subcodes of the binary second-order Reed--Muller code \(RM_2(2,n)\). We determine
several exact all-parameter regions of this spectrum. First, for every \(n\ge4\),
the minimum value \(7\cdot2^{n-4}\) is isolated and the next possible support is
at least \(2^{n-1}\). Second, for every \(n\ge7\), a nonfull support is at most
\(2^n-2^{n-6}\), and this bound is sharp. We determine \(\mathcal S_n^{(3)}\)
exactly for \(2\le n\le9\), and for every \(n\ge8\) we classify the complete
near-full interval corresponding to a nonempty common zero fibre smaller than
\(2^{n-5}\). We also prove that the support
\(493\cdot2^{n-9}\) is impossible for every \(n\ge12\).

The structural part of the paper gives a finite sound-and-complete recursion
whenever one nonzero output combination has nonzero zero-frequency Walsh
coefficient and polar rank two or four. The recursion uses the canonical
decomposition of a binary alternating pencil, exact affine mask and sign
transitions for all Kronecker, degenerate regular, and nondegenerate regular
blocks, and an explicit independence filter. It is constructive: every accepted
state can be traced back to a three-dimensional quadratic subcode. We do not
determine the complete third support spectrum; the Walsh-zero low-rank layer and
the layer in which every nonzero Walsh component has polar rank at least six
remain outside the theorem.

\end{abstract}

\section{Introduction}


For a linear code \(C\subseteq\mathbb F_2^N\), the support of a subcode
\(D\le C\) is the union of the supports of its codewords. The generalized
Hamming weights retain only the smallest support in each subcode dimension.
The corresponding higher support spectrum asks for every attainable support
size, and is usually much less understood. Even for Reed--Muller codes, complete
higher spectra are known only in restricted parameter ranges; see Ghorpade,
Johnsen, Ludhani, and Pratihar \cite{ghorpade2026} for recent results and a survey of this
distinction.

We study three-dimensional subcodes of

\[
RM_2(2,n)=\{\text{Boolean functions on }\mathbb F_2^n
\text{ of degree at most two}\}.
\]

Write

\[
\mathcal S_n^{(3)}=
\{|\mathrm{Supp}(D)|:D\le RM_2(2,n),\ \dim D=3\}.
\]

The analogous two-dimensional support-position problem has an all-dimensional
classification \cite{carptopus2026}. Passing from a quadratic pencil to a three-parameter
quadratic net changes the structure substantially. Seven nonzero codewords must
be compatible on a Fano parameter plane, while their affine Walsh data come
from one common three-dimensional system. This prevents the ordinary quadratic
weight spectrum from being used independently in each of the seven directions.

The present paper does not solve the complete three-dimensional spectrum.
Instead, it closes five natural and mutually reinforcing parts of it.

\begin{enumerate}
\item Two sharp all-dimensional gaps are determined at the minimum and at full
support.
\item The exact spectra are determined for \(2\le n\le9\).
\item The entire near-full interval below twice the degree-six minimum weight is
classified for every \(n\ge8\).
\item A new infinite forbidden ray is proved:
\(493\cdot2^{n-9}\notin\mathcal S_n^{(3)}\) for \(n\ge12\).
\item All systems having a nonzero-Walsh direction of polar rank two or four are
described by a finite complete canonical-pencil recursion.
\end{enumerate}


The first three parts combine classical Reed--Muller weight restrictions with
Fano-plane rank geometry and explicit slice constructions. The last two require
affine information on alternating pencils, rather than only their polar ranks.
The canonical classification itself is classical \cite{plut2015,macariorat2013}. Our contribution is an
exact quotient of every canonical block into the masks, sign functions, and
relation kernels needed for a common zero fibre, valid for arbitrary block
length and compatible with constructive dynamic programming.

The proofs mix ordinary arguments with finite exact certificates. We separate
their roles explicitly in Section 8. In particular, a finite computation is
never used to extrapolate a theorem to unbounded \(n\).

\section{Quadratic systems, support, and Walsh data}


Let

\[
F=(f_1,f_2,f_3):\mathbb F_2^n\longrightarrow\mathbb F_2^3
\]

have linearly independent components of degree at most two, and put

\[
Z(F)=|F^{-1}(0)|.
\]

The corresponding subcode \(D=\langle f_1,f_2,f_3\rangle\) satisfies

\[
|\mathrm{Supp}(D)|=2^n-Z(F).                    \tag{2.1}
\]

For \(a\in\mathbb F_2^3\setminus\{0\}\), set \(f_a=a\cdot F\), and write

\[
W_a=\sum_{x\in\mathbb F_2^n}(-1)^{f_a(x)}.
\]

Pointwise, each nonzero value of \(F(x)\) is detected by exactly four of the
seven functions \(f_a\). Consequently

\[
|\mathrm{Supp}(D)|
=\frac14\sum_{a\ne0}\mathrm{wt}(f_a)
=7\,2^{n-3}-\frac18\sum_{a\ne0}W_a,             \tag{2.2}
\]

and Fourier inversion gives

\[
8Z(F)=2^n+\sum_{a\ne0}W_a.                      \tag{2.3}
\]

The polar form of a quadratic function \(f\) is

\[
B_f(x,y)=f(x+y)+f(x)+f(y)+f(0).
\]

It is alternating. If its rank is \(r\), then \(r\) is even and

\[
W_f(0)\in\{0,\pm2^{n-r/2}\}.                    \tag{2.4}
\]

The coefficient is zero precisely when the affine linear part does not vanish
on the radical. If it is nonzero, its sign is the value of the induced
nondegenerate quadratic form after completing the square.

Equation (2.4) yields the divisibility

\[
2^{\max(0,\lceil(n-6)/2\rceil)}\mid Z(F).        \tag{2.5}
\]

Indeed, in even dimension \(n=2m\), every nonzero Walsh coefficient is divisible
by \(2^m\); in odd dimension \(n=2m+1\), it is divisible by \(2^{m+1}\). Apply
this to (2.3), including the small-dimensional boundary separately.

We use two classical facts about Reed--Muller weights.

\Needspace{0.23\textheight}
\begin{fact}[minimum and next weight]
A nonzero Boolean polynomial of degree at
most \(r\) has weight at least \(2^{n-r}\). For \(2\le r\le n-2\), a weight
strictly between \(2^{n-r}\) and \(3\cdot2^{n-r-1}\) does not occur.
\end{fact}


\Needspace{0.23\textheight}
\begin{fact}[weights below twice the minimum]
If
\(p\in RM_2(r,n)\) has
\end{fact}


\[
0<\mathrm{wt}(p)<2^{n-r+1},
\]

then

\[
\mathrm{wt}(p)=2^{n-r+1}-2^i
\]

for some integer \(i\). This is the low-weight theorem of Berlekamp and Sloane
\cite{berlekamp1969}, with its stated boundary included.

For the upper edge we use Warning's second theorem in its fibre form: a nonempty
common zero set of polynomials whose degree sum is \(d<n\) has at least
\(2^{n-d}\) points \cite{asgarli2018}.

\section{Two all-dimensional gaps and the exact spectra through nine variables}


\subsection{The isolated minimum}


\Needspace{0.23\textheight}
\begin{theorem}
For every \(n\ge4\) and every three-dimensional
\(D\le RM_2(2,n)\),
\end{theorem}


\[
|\mathrm{Supp}(D)|=7\cdot2^{n-4}
\quad\text{or}\quad
|\mathrm{Supp}(D)|\ge2^{n-1}.                  \tag{3.1}
\]

Both values are attained.

\begin{proof}
A nonzero word of \(RM_2(2,n)\) has minimum weight \(2^{n-2}\), and
every nonminimum nonzero word has weight at least \(3\cdot2^{n-3}\). If at most
five of the seven nonzero words in \(D\) have minimum weight, their average in
(2.2) already gives \(|\mathrm{Supp}(D)|\ge2^{n-1}\).

Suppose that six words have minimum weight. Their six polar matrices have rank
two. Every \(4\times4\) Pfaffian of the three-parameter polar net restricts to a
homogeneous quadratic function on \(\mathbb F_2^3\). Such a function cannot be
nonzero at exactly one of the seven nonzero parameter points. Hence all
Pfaffians also vanish at the seventh point, whose polar rank is therefore at
most two. If the seventh word is not minimum, its weight is at least
\(2^{n-1}\), and the same average again gives the second alternative in (3.1).
Otherwise all seven words are minimum and equality holds in the first
alternative.

The subcode \(x_1\langle x_2,x_3,x_4\rangle\) attains the minimum. The subcode
\(\langle x_1,x_1x_2,x_1x_3\rangle\) attains \(2^{n-1}\), after adjoining dummy
variables as needed.
\end{proof}


\subsection{The last nonfull support}


\Needspace{0.23\textheight}
\begin{theorem}
For every \(n\ge7\), if
\(D\le RM_2(2,n)\) is three-dimensional and does not have full support, then
\end{theorem}


\[
|\mathrm{Supp}(D)|\le2^n-2^{n-6}.               \tag{3.2}
\]

The bound and full support are both attained.

\begin{proof}
The complement of the support is the common zero set of three
quadratic functions. If nonempty, Warning's second theorem gives at least
\(2^{n-6}\) points. For equality, on six distinct variables take

\[
f_1=x_1x_2+x_1+x_2,\quad
f_2=x_3x_4+x_3+x_4,\quad
f_3=x_5x_6+x_5+x_6.
\]

Their common zero set fixes the first six variables to zero and has size
\(2^{n-6}\). A subcode containing the constant function one has full support.
\end{proof}


\subsection{Exact small spectra}


We first record the rank geometry used by the finite arithmetic certificate.
Identify an alternating form \(A\) with the linear map \(V\to V^*\) that it
induces, and put \(\operatorname{Ess}(A)=\operatorname{im}A\). For rank two this
is the two-dimensional essential support; equivalently, if \(A=u\wedge v\),
then \(\operatorname{Ess}(A)=\langle u,v\rangle\).

\Needspace{0.23\textheight}
\begin{lemma}[Fano rank rules]
Let
\(T:\mathbb F_2^3\to\operatorname{Alt}(V)\) be linear, and evaluate the ranks of
\(T\) at the seven nonzero parameter points.
\end{lemma}


\begin{enumerate}
\item If \(T\) is not injective, one of the seven images is zero and, after one
zero image is removed, every rank occurs with even multiplicity.
\item If at least four parameter points have rank two, then every image of \(T\)
has rank at most four.
\item Suppose exactly three parameter points have rank two and some image has
rank eight. Then \(T\) is injective, the three rank-two points form a line
\(\{A,B,A+B\}\), and all four points outside that line have rank at least
six. If \(\dim V\in\{8,9\}\), at least two of those four points have rank
eight.
\item Suppose exactly two parameter points \(A,B\) have rank two, \(T\) is
injective, and an outside point \(C\) has rank eight. Then \(A+B\) has rank
four and \(C,C+A,C+B\) have rank at least six.
\item If six parameter points have rank two, then the seventh has rank at most
two.
\end{enumerate}


\begin{proof}
If \(T\) has a kernel of dimension one, the three nonzero elements of
its two-dimensional image each have two preimages among the seven nonzero
parameters, while the nonzero kernel point maps to zero. If the kernel has
dimension two, the remaining zero rank has multiplicity two and the unique
nonzero image rank has multiplicity four; if \(T=0\), the remaining zero rank
has multiplicity six. This proves part 1.

For part 2, first suppose that \(T\) is not injective. Its image has dimension
at most two; image dimension at most one is immediate, while in dimension two
part 1 forces two independent nonzero image forms to have rank two. Their sum
has rank at most four, so all images have rank at most four. Now suppose that
\(T\) is injective. Choose
three independent rank-two forms \(A,B,C\) among four such points. Relative to
this basis the fourth is either a pairwise sum or \(A+B+C\). In the first case,
say \(A+B\) has rank two. Rank-two alternating forms are decomposable, and the
sum of two such forms is decomposable only when their essential two-planes
meet. Thus the essential supports of \(A,B,C\) span at most five dimensions.
In the second case, if the three essential supports were in direct sum then
\(A+B+C\) would have rank six, so again their span has dimension at most five.
Every image of \(T\) is supported there and therefore has even rank at most
four.

For part 3, part 1 makes \(T\) injective. If the three rank-two points were not
collinear, they would form a basis and all image forms would be supported on at
most six dimensions, contradicting the rank-eight point. Hence they are
\(A,B,A+B\). Rank subadditivity gives
\(\operatorname{rank}(C+L)\ge6\) for every point \(L\) on this line. When
\(\dim V\) is eight or nine, let \(P(X)\) denote respectively the maximal
Pfaffian or the vector of maximal principal Pfaffians. The three decomposable
forms on the line have a common wedge factor, so the matching expansion has no
repeated-update or mixed-update terms and gives

\[
P(C+A)+P(C+B)+P(C+A+B)=P(C).
\]

The right-hand side is nonzero. At least one updated form therefore also has
rank eight, proving the last assertion.

For part 4, injectivity and the assumption that there are exactly two
rank-two points show that \(A+B\) is nonzero and not rank two; rank
subadditivity makes its rank at most four, hence exactly four. The same
inequality applied to \(C=(C+A)+A=(C+B)+B\) gives the two rank-six lower
bounds.

Finally, restrict any four-dimensional principal Pfaffian to the Fano parameter
plane. It is a homogeneous quadratic Boolean function. Such a function cannot
vanish at six nonzero points and be nonzero only at the seventh: its truth-set
size is even. Thus every four-dimensional Pfaffian also vanishes at the seventh
point, proving part 5.
\end{proof}


\Needspace{0.23\textheight}
\begin{theorem}
The exact support-position sets for \(2\le n\le9\) are
\end{theorem}


\[
\begin{aligned}
\mathcal S_2^{(3)}&=\{3,4\},\\
\mathcal S_3^{(3)}&=\{4,5,6,7,8\},\\
\mathcal S_4^{(3)}&=\{7,8,\ldots,16\},\\
\mathcal S_5^{(3)}&=\{14\}\cup\{16,17,\ldots,32\},\\
\mathcal S_6^{(3)}&=\{28,32,34\}\cup\{36,37,\ldots,64\},\\
\mathcal S_7^{(3)}&=\{56,64,68\}\cup\{72,74,\ldots,128\},
\end{aligned}                                                    \tag{3.3}
\]

where the last progression contains even values only, and

\[
\begin{aligned}
\mathcal S_8^{(3)}={}&
\{112,128,136,144,148,152,156,160,164\}\\
&\cup\{166,168,\ldots,252\}\cup\{256\},\\
\mathcal S_9^{(3)}={}&2\mathcal S_8^{(3)}.
\end{aligned}                                                    \tag{3.4}
\]

Again, the progression in \(\mathcal S_8^{(3)}\) has step two.

We give the proof in a form that separates general exclusions from finite
positive certificates.

For \(n=2\), a three-dimensional subcode is a hyperplane in
\(\mathbb F_2^4\), so its support is three or four. For \(n=3\), the code is the
even-weight code of length eight. A full-support three-dimensional even code
of each length \(4,\ldots,8\) is obtained from nonzero columns of
\(\mathbb F_2^3\) with zero column sum and full span. For \(n=4\), Theorem 3.1
leaves the isolated value seven and the interval from eight to sixteen; a
fixed four-variable polar net realizes every affine-orbit support in that
interval.

For the higher small dimensions we use the following slice identity. If
\(F,L:\mathbb F_2^m\to\mathbb F_2^3\), with \(F\) quadratic and \(L\) affine,
then

\[
H(x,t)=F(x)+tL(x)
\quad\Longrightarrow\quad
|\mathrm{Supp}(H)|=|\mathrm{Supp}(F)|+|\mathrm{Supp}(F+L)|.       \tag{3.5}
\]

Three fixed polar nets have affine-orbit support sets

\[
\begin{aligned}
R_4&=\{8,9,\ldots,16\},\\
R_5&=\{17\}\cup\{19,20,\ldots,32\},\\
R_7&=\{78,82,86\}\cup\{88,90,\ldots,124,128\}.
\end{aligned}                                                    \tag{3.6}
\]

Their three polar components are linearly independent. Explicit quadratic
parts for these nets and an exhaustive enumeration of all vector-valued affine
offsets are included in the verification package. Thus
\(R_4+R_4=[16,32]\),
\(R_5+R_5=\{34\}\cup[36,64]\), and the sumset of \(R_7\), together with
doubling from dimension seven, realizes every retained value in (3.4).

It remains to justify the exclusions. In dimensions six and seven, the only
arithmetic candidates not constructed are respectively \(33,35\) and
\(66,70\). Each would require at least four rank-two points together with a
polar form of rank at least six in the same Fano parameter plane. This is
impossible: if the polar map is noninjective, the image has dimension at most
two and the third nonzero image has rank at most four; if it is injective, four
rank-two points force the combined support of a basis of rank-two forms into at
most five dimensions.

In dimensions eight and nine, every seven-word weight multiset for an
unconstructed arithmetic candidate is enumerated, with every balanced word
branched over its possible even polar ranks. Lemma 3.3 reduces the
candidate lists to zero except for the apparent nine-dimensional pattern

\[
(2,2,2,6,6,6,8).                                \tag{3.7}
\]

The three rank-two points must be collinear. If \(A,B,A+B\) are the forms on
that line and \(C\) has rank eight, the vector of maximal Pfaffians satisfies

\[
P(C+A)+P(C+B)+P(C+A+B)=P(C).                    \tag{3.8}
\]

This holds in dimensions eight and nine: the rank-two line has a common wedge
factor, so the Pfaffian expansion has neither repeated-update nor mixed-update
terms. Since \(P(C)\ne0\), at least one updated form also has rank eight,
contradicting (3.7). Finally, the apparent supports \(254\) and \(508\) would
leave only two or four common zeros, respectively, contrary to Theorem 3.2.
The verification program exhausts all seven-word quadratic-weight multisets,
branches over every even polar rank compatible with a balanced word, and then
applies only Lemma 3.3 and (3.8). Its finite output is therefore an arithmetic
certificate rather than an unproved rank-geometry oracle. This proves
(3.3)--(3.4).
\hfill$\square$


\section{The exact near-full interval}


The common-zero indicator is

\[
\mathbf1_{F=0}=(1+f_1)(1+f_2)(1+f_3),           \tag{4.1}
\]

a Boolean polynomial of degree at most six.

\Needspace{0.23\textheight}
\begin{theorem}
Put \(M=2^{n-6}\). For every \(n\ge8\),
\end{theorem}


\[
0<Z(F)<2M
\]

holds if and only if

\[
Z(F)=2M-2^i,
\qquad
\left\lceil\frac{n-6}{2}\right\rceil\le i\le n-6.              \tag{4.2}
\]

Equivalently, all supports strictly between \(2^n-2^{n-5}\) and \(2^n\) are

\[
2^n-\left(2^{n-5}-2^{n-\mu-5}\right),
\qquad
1\le\mu\le\left\lfloor\frac{n-4}{2}\right\rfloor.             \tag{4.3}
\]

\begin{proof}
Apply Fact 2.2 to (4.1). It gives
\(Z(F)=2M-2^i\). Divisibility (2.5) forces
\(i\ge\lceil(n-6)/2\rceil\), while positivity and the strict upper bound force
\(i\le n-6\).

For sufficiency, use disjoint variable blocks and set

\[
g_1=x_1x_2,\qquad g_2=x_3x_4,
\qquad
g_3=y_1y_2+\cdots+y_{2\mu-1}y_{2\mu}.
\]

The third form is nondegenerate of positive type and has
\(2^{2\mu-1}-2^{\mu-1}\) ones. Take \(f_j=1+g_j\) and add dummy variables.
The three polar parts lie on disjoint blocks and are linearly independent, and

\[
Z(F)=\mathrm{wt}(g_1g_2g_3)
=2^{n-5}-2^{n-\mu-5}.
\]

The substitution \(i=n-\mu-5\) gives precisely the range in (4.2).
\end{proof}


The strict boundary is important. The value \(Z=2M\) is attained by systems
having no minimum quadratic factor and is not classified by Theorem 4.1.

\section{Arithmetic forcing for the \texorpdfstring{$19/8$}{19/8} ray}


We now study

\[
Z(F)=19\cdot2^{n-9}.                             \tag{5.1}
\]

The proof first forces a low-polar-rank negative Walsh direction, and then uses
the complete recursions of Section 6.

\subsection{The rank-two case}


Suppose a negative-Walsh rank-two direction occurs. After an output change of
basis, its indicator factor has minimum weight and may be normalized to
\(x_1x_2\). Restricting to its codimension-two affine support reduces (5.1) to
the intersection of the one-sets of two quadratic functions \(a,b\) in
\(m=n-2\) variables. Fourier inversion gives

\[
4|\{a=b=1\}|=2^m-W_a-W_b+W_{a+b}.               \tag{5.2}
\]

At the target (5.1), division by \(2^{m-5}\) asks for three signed powers of two
whose sum is \(13\). Negative exponents cannot occur. Indeed, at the least
negative exponent exactly two of the three nonzero terms occur; opposite signs
cancel and leave the impossible single power \(13\), while equal signs combine
one level higher and either demand \(12\) or \(14\), or reduce to zero or a
single signed power. Among integer powers, the only possibilities, up to order,
are

\[
16-2-1,\qquad 16-4+1,\qquad 8+4+1.             \tag{5.3}
\]

Their polar-rank multisets are \(\{2,8,10\}\), \(\{2,6,10\}\), and
\(\{4,6,10\}\). The middle case violates rank subadditivity. In each remaining
case the largest rank equals the sum of the other two. Their polar images are
therefore disjoint, and Fourier convolution at zero gives

\[
W_{a+b}(0)=2^{-m}W_a(0)W_b(0).                  \tag{5.4}
\]

The signs required by (5.3) have negative product, contradicting (5.4). Thus
the target cannot contain a negative-Walsh rank-two direction.

\subsection{Forcing rank two or four}


Normalize the seven Walsh coefficients by

\[
c_a=\frac{W_a}{2^{n-6}}.
\]

Equation (2.3) and (5.1) give \(\sum_{a\ne0}c_a=-45\), where every nonzero
coefficient is \(\pm2^{6-r_a/2}\). The endpoints \(+64\) and \(-64\) are
impossible: the former makes an output combination zero, and the latter makes
it one and hence forces \(Z=0\).

Assume neither \(-32\) nor \(-16\) occurs. Nonzero fractional terms cannot
occur: one fractional term cannot give an integer sum, while \(k\ge2\) such
terms give a total strictly greater than

\[
-k-8(7-k)=-56+7k\ge-42.
\]

Positive \(16\) or \(32\) is also impossible because the remaining six terms,
even if all \(-8\), do not reach \(-45\). Starting from seven copies of \(-8\),
the required increase is \(11\). The only admissible decomposition is
\(11=4+7\), so the coefficient multiset must be

\[
(-8,-8,-8,-8,-8,-4,-1).                        \tag{5.5}
\]

Through the rank-twelve point, two of the three Fano lines have rank multiset
\((12,6,6)\). Equality in rank subadditivity again forces the Walsh sign of the
rank-twelve member to be the product of the two rank-six signs, contrary to all
three signs being negative. We have proved:

\Needspace{0.23\textheight}
\begin{lemma}
Every system satisfying (5.1) has a negative-Walsh direction of
polar rank two or four.
\end{lemma}


Section 6 gives a complete finite recurrence for both ranks. The rank-two case
has already been excluded above. The rank-four recurrence excludes every
remaining state and therefore closes the ray in Section 7.

\section{A finite complete recursion for nonzero-Walsh polar rank two or four}


\subsection{Core reduction}


Choose a nonzero output combination \(f_0\) with \(W_{f_0}(0)\ne0\) and polar
rank \(2\rho\), where \(\rho\in\{1,2\}\). After an affine change of variables,
the factor \(q_0=1+f_0\) is the pullback from \(U\) of one of the two
nonsingular quadrics

\[
T\subseteq U=\mathbb F_2^{2\rho},
\qquad
|T|=2^{2\rho-1}\pm2^{\rho-1}.                  \tag{6.1}
\]

Thus \(|T|\) is one or three for \(\rho=1\), and six or ten for \(\rho=2\).
Write the ambient space as \(U\oplus R\), \(m=\dim R\). The restrictions to
\(R\) of the other two polar forms constitute an alternating pencil. Their
cross terms with \(U\), affine parts, and constants vary with \(u\in U\).

For each canonical block of the pencil, and for each of the three nonzero
combinations of the two residual quadratic functions, record:

\begin{itemize}
\item a mask \(M_i\subseteq U\) on which the Walsh coefficient is nonzero;
\item an exponent \(e_i\) giving its absolute value \(2^{e_i}\);
\item a set \(N\subseteq M_1\cap M_2\cap M_3\) on which the product of the three
nonzero Walsh signs is negative;
\item a relation kernel \(\Lambda\le\mathbb F_2^2\) recording linear relations
between the two noncore residual functions.
\end{itemize}


A state is therefore

\[
\Sigma=(m;e_1,e_2,e_3;M_1,M_2,M_3;N;\Lambda).  \tag{6.2}
\]

Orthogonal direct sum gives the exact composition law

\[
\begin{aligned}
m''&=m+m',& e_i''&=e_i+e_i',\\
M_i''&=M_i\cap M_i',&
N''&=(N\mathbin\triangle N')\cap M_1''\cap M_2''\cap M_3'',\\
\Lambda''&=\Lambda\cap\Lambda'.
\end{aligned}                                                     \tag{6.3}
\]

\subsection{Canonical blocks and affine-product classes}


Let \(\mathcal A_U\) be the truth masks of affine functions on \(U\), and let
\(\mathcal Q_2(U)\) denote the Boolean functions on \(U\) of algebraic degree
at most two. Also put

\[
\mathcal P_r(U)=
\left\{\mathop{\triangle}_{j=1}^r(A_j\cap B_j):
A_j,B_j\in\mathcal A_U\right\},
\qquad \mathcal P_0(U)=\{\varnothing\}.         \tag{6.4}
\]

The alternating-pencil canonical decomposition consists of the following
blocks \cite{plut2015,macariorat2013}.

\textbf{Singular Kronecker block \(K_k\).} In dimension \(2k+1\),

\[
Q_1=\sum_{i=0}^{k-1}x_iy_i,
\qquad
Q_2=\sum_{i=0}^{k-1}x_{i+1}y_i.                \tag{6.5}
\]

For \(k\ge1\), the three masks independently range over \(\mathcal A_U\), all
exponents are \(k+1\), and, on their common mask \(C\),

\[
N=P\cap C,\qquad P\in\mathcal P_{k-1}(U).      \tag{6.6}
\]

The two residual polar generators are independent. The block \(K_0\) is a
one-dimensional common radical. For affine zero masks \(A,B\), its masks are

\[
(A,B,U\setminus(A\mathbin\triangle B)),
\]

its product sign is positive, and its relation kernel is read directly from
the two affine coefficients.

\textbf{Degenerate regular chain \(D_k^{(j)}\).} In dimension \(2k\), one of the
three pencil points is degenerate. Its mask lies in \(\mathcal P_1(U)\), the
other two masks are all of \(U\), its exponent is \(k+1\), and the other two
exponents are \(k\). On the common mask, (6.6) again holds. Only the length-one
case with an identically zero degenerate residual function contributes a
nontrivial relation line.

\textbf{Nondegenerate regular part \(R_d\).} Every component on which the three
rational pencil points are nondegenerate has, after local coordinates,

\[
Q_1=x^{\mathsf T}y,
\qquad Q_2=x^{\mathsf T}Cy,
\qquad C,I+C\text{ invertible}.                 \tag{6.7}
\]

All three masks are \(U\), all exponents are \(d\), and
\(N\in\mathcal P_d(U)\). This transfer depends only on the total half-dimension
\(d\), not on the elementary divisors of \(C\).

The last assertion follows from the exact factorization of the product-sign
control matrix. Put \(D=C^{\mathsf T}\) and \(A=(I+D)^{-1}\). Then

\[
M_C=
\begin{pmatrix}I+A&A\\A&D^{-1}+A\end{pmatrix}
=\begin{pmatrix}D\\I\end{pmatrix}
A
\begin{pmatrix}I&D^{-1}\end{pmatrix},           \tag{6.8}
\]

which has rank \(d\). The analogous control ranks for \(K_k\) and \(D_k\) are
both \(k-1\). Explicit triangular elimination gives, for all block lengths,

\[
P_K=\sum_{i=0}^{k-2}(p_i+s_{i+1})
\left(\sum_{j=i+1}^{k-1}r_j+\sum_{j=i}^{k-2}t_j\right),          \tag{6.9}
\]

and

\[
P_D=\sum_{i=1}^{k-1}\bigl((p_1)_i+(p_2)_{i-1}\bigr)
\left(\sum_{j=0}^{i-1}(r_1)_j+\sum_{j=1}^{i}(r_2)_j\right).     \tag{6.10}
\]

The left and right factors in each sum can be specified independently of the
nonzero masks. Equations (6.8)--(6.10), rather than a bounded block table, prove
the arbitrary-length transition formulas.

\subsection{Terminal output and independence}


For each \(\rho\in\{1,2\}\) and each of the two choices of \(T\), set
\(m=n-2\rho\). Starting with the empty state, append every canonical block
transition whose accumulated dimension does not exceed \(m\), using (6.3).
Retain only block sequences whose accumulated dimension is exactly \(m\);
these and only these states are terminal. Take the union of their outputs.
There are finitely many states for fixed \(m\), because every block parameter
is bounded by its dimension.

For a terminal state and \(u\in U\), let

\[
a_i(u)=
\begin{cases}
2^{e_i},&u\in M_i,\\
0,&u\notin M_i,
\end{cases}
\]

and reverse the sign of \(a_3(u)\) for \(u\in N\). If the two pure core
quadratic functions take values \(c_1(u),c_2(u)\), the residual pair-intersection
number is

\[
L_u=\frac{2^m-(-1)^{c_1(u)}a_1(u)-(-1)^{c_2(u)}a_2(u)
+(-1)^{c_1(u)+c_2(u)}a_3(u)}4.                  \tag{6.11}
\]

Enumerate the two actual functions \(c_1,c_2\in\mathcal Q_2(U)\) and sum (6.11) over
\(u\in T\). The resulting system is three-dimensional precisely when, for every
nonzero \((b_1,b_2)\in\Lambda\),

\[
b_1f_1^U+b_2f_2^U\notin\{0,f_0\}.              \tag{6.12}
\]

This condition is essential on the three-point and ten-point quadrics; removing
it admits false outputs coming from dependent generators.

\Needspace{0.23\textheight}
\begin{theorem}[finite low-rank recursion]
Fix \(n\). Among all
three-dimensional systems \(F\) having a nonzero output combination with
nonzero zero-frequency Walsh coefficient and polar rank two or four, the
attainable values of \(Z(F)\) are exactly the values produced by the finite
recursion (6.2)--(6.12). Each accepted value comes with a reconstructible
quadratic witness.
\end{theorem}


\begin{proof}
For necessity, normalize the selected direction and decompose the
restricted alternating pencil on its radical into canonical \(K_k\),
\(D_k^{(j)}\), and regular primary components. Equations (6.6)--(6.10) give the
exact transition of every block; (6.3) is Walsh multiplication under orthogonal
sum; and (6.11) is Fourier inversion on each core slice. Hence every system
reaches a terminal state and output accepted by (6.12).

For sufficiency, every block transition has explicit affine coefficients.
Control-matrix factorization reconstructs the prescribed sign mask, and
orthogonal sum reconstructs the terminal state. The chosen core functions
produce (6.11), while (6.12) removes exactly the possible generator relations.
Parent pointers, block parameters, affine-product factors, and core functions
therefore reconstruct a linearly independent triple realizing every accepted
value.
\end{proof}


The theorem deliberately requires a \emph{nonzero} Walsh coefficient. There are
infinite systems with rank profile \((2,2,0,2,0,0,2)\) for which all seven
zero-frequency Walsh coefficients vanish. They are a boundary control, not an
input to Theorem 6.1.

\section{An infinite forbidden support ray}


\Needspace{0.23\textheight}
\begin{theorem}
For every \(n\ge12\),
\end{theorem}


\[
19\cdot2^{n-9}\notin\{Z(F):F\text{ has three independent quadratic components}\}.
\]

Equivalently,

\[
493\cdot2^{n-9}\notin\mathcal S_n^{(3)}.         \tag{7.1}
\]

\begin{proof}
Lemma 5.1 forces a negative-Walsh direction of polar rank two or four.
Section 5.1 excludes the rank-two case.

For rank four, specialize Theorem 6.1. If \(m=n-4\) and
\(\delta_i\) are the three residual pencil half-ranks, the normalized target
equation is

\[
\sum_{i=1}^3 c_i2^{5-\delta_i}=-116,
\qquad
c_i\in\{-6,-4,-3,-2,-1,0,1,2,3,4,6\}.          \tag{7.2}
\]

Together with the half-rank triangle inequalities, (7.2) leaves only

\[
(0,3,3),(0,4,4),(0,5,5),(1,2,3),
(1,3,3),(1,3,4),(1,4,4).                        \tag{7.3}
\]

Thus, apart from \(K_0\), only \(K_1,D_1,D_2\) can occur. Their integer
decompositions give nine bounded canonical cores. Multiple \(K_0\) layers
stabilize after two blocks under the intersection law (6.3), so it is enough to
check zero, one, and at least two common-radical blocks for each core. Substituting
the exact masks and sign classes from Section 6 into (6.11) excludes (7.2) in
every case. This is a finite exact consequence of the arbitrary-length block
theorem, not a scan over bounded \(n\).
\end{proof}


At \(n=12\), an independent enumeration of all 65 residual pencil exponent
profiles and 27 surviving canonical decompositions gives the same exclusion.
The only state surviving an artificial independent-sign relaxation requires a
five-point negative mask omitted by the exact length-two \(D_2\) transfer. This
provides a useful destructive control for the sign coupling in the proof.

\section{Computational proof boundary and reproducibility}


The main results use finite exact certificates at several explicitly identified
interfaces.

\subsection{General arguments independent of finite enumeration}


\begin{itemize}
\item the support/Walsh identities (2.1)--(2.3) and divisibility (2.5);
\item Theorems 3.1 and 3.2;
\item the explicit constructions in Theorem 4.1;
\item the signed-dyadic and rank-equality arguments in Section 5;
\item canonical-pencil completeness and the arbitrary-length transition formulas
(6.8)--(6.10);
\item termination, soundness, completeness, and witness reconstruction in
Theorem 6.1;
\item the dimension-independent reduction (7.2)--(7.3).
\end{itemize}


\subsection{Finite exact certificates}


\begin{itemize}
\item affine-orbit support sets of the three explicit polar nets in (3.6), their
polar-component independence, and the corresponding slice sumsets;
\item all seven-word weight and polar-rank patterns used for the \(n=8,9\)
exclusions, including noninjective polar maps and the maximal-Pfaffian
negative control;
\item exact formula checks for \(K_k,D_k,R_d\), tested on representative block
lengths and regular companion matrices;
\item finite state composition, relation-kernel intersection, one/three/six/ten-point
terminal output, and independence positive and negative controls;
\item the nine rank-four cores and stabilized \(K_0\) layers in Theorem 7.1.
\end{itemize}


The verification runner recomputes every retained certificate using exact
integer, rational, or \(\mathbb F_2\) arithmetic. One fixed-seed routine samples
regular matrices only as a regression for the symbolic factorization (6.8);
the sample is not used to infer generality. A failed search and a solver timeout
are never used as nonexistence evidence.

\section{Relation to prior work and limits}


Ghorpade, Johnsen, Ludhani, and Pratihar \cite{ghorpade2026} determine higher support spectra for
\(RM_q(2,2)\), where the number of variables is fixed and the field varies. The
present problem fixes \(q=2\), lets the number of variables vary, and studies
three-dimensional subcodes. The results are adjacent but neither specializes
to the other.

Berlekamp and Sloane \cite{berlekamp1969} provide the ordinary low-weight restriction used as a
necessary input in Theorem 4.1. Ordinary codeword weights do not supply the
compatibility or sufficiency of a common three-quadratic zero fibre. The latter
is provided here by explicit independent constructions.

The alternating-pencil decompositions of Plût, Fouque, and Macario-Rat \cite{plut2015} and
Macario-Rat, Plût, and Gilbert \cite{macariorat2013} are structural inputs to Theorem 6.1. We do
not claim a new canonical classification. The new part is the exact affine
Walsh quotient of each block, the arbitrary-length product-sign formulas, and
the complete common-zero recursion with independence filtering.

The previously published two-dimensional support-position classification \cite{carptopus2026}
is a precursor, not a proof of the present results. A three-dimensional system
contains seven coupled directions and requires the Fano and canonical-pencil
mechanisms developed here.

The following questions remain outside the scope of this paper.

\begin{enumerate}
\item We do not determine the complete set \(\mathcal S_n^{(3)}\) for every \(n\).
\item The low-polar-rank recursion excludes Walsh-zero selected directions; the
all-high-rank layer is also open.
\item We do not count subcodes of a given support size or classify quadratic nets
up to equivalence.
\item We make no claim over odd characteristic or arbitrary finite fields.
\item The current literature search found no direct coverage of the frozen
theorems, but it does not establish global priority or external acceptance.
\end{enumerate}


\section{Conclusion}


The third support spectrum of \(RM_2(2,n)\) is not a free seven-fold combination
of ordinary quadratic weights. Nevertheless, several natural regions admit
exact structure. The two spectrum edges have sharp universal gaps; all spectra
through nine variables can be closed; the complete near-full interval follows
from a classical low-weight theorem plus a compatible three-factor construction;
and the \(19/8\) common-zero ray is uniformly forbidden.

The canonical-pencil recursion explains why low polar rank remains tractable.
After fixing one rank-two or rank-four direction, the remaining wild-looking net
becomes an alternating pencil on the radical. Full marked-module orbits are not
needed: affine masks, product-sign classes, and a five-valued relation kernel
form a finite sufficient quotient. The unresolved part is therefore sharply
located in the Walsh-zero low-rank boundary and in nets whose every nonzero
Walsh direction has polar rank at least six.

\begin{thebibliography}{9}
\footnotesize
\raggedright

\bibitem{carptopus2026}
Carptopus,
\emph{Attainable second support weights of binary second-order Reed--Muller codes},
MathExplore Notes, Public Beta v0.2 (2026).
\href{https://doi.org/10.5281/zenodo.22059063}{doi:10.5281/zenodo.22059063}.

\bibitem{ghorpade2026}
S.~R. Ghorpade, T. Johnsen, R. Ludhani, and R. Pratihar,
\emph{Betti Numbers and Higher Weight Spectra of Reed--Muller Codes $RM_q(2,2)$},
arXiv:2408.02548v3 (2026).
\href{https://arxiv.org/abs/2408.02548v3}{arXiv:2408.02548v3}.

\bibitem{berlekamp1969}
E.~R. Berlekamp and N.~J.~A. Sloane,
\emph{Restrictions on Weight Distribution of Reed--Muller Codes},
Information and Control 14 (1969), 442--456.
\href{https://doi.org/10.1016/S0019-9958(69)90150-8}{doi:10.1016/S0019-9958(69)90150-8}.

\bibitem{asgarli2018}
S. Asgarli,
\emph{A new proof of Warning's second theorem},
American Mathematical Monthly 125 (2018), 548--552.
\href{https://doi.org/10.1080/00029890.2018.1449555}{doi:10.1080/00029890.2018.1449555}.

\bibitem{plut2015}
J. Pl\^ut, P.-A. Fouque, and G. Macario-Rat,
\emph{Solving the Isomorphism of Polynomials with Two Secrets Problem for all Pairs of Quadratic Forms},
arXiv:1406.3163v3 (2015).
\href{https://arxiv.org/abs/1406.3163v3}{arXiv:1406.3163v3}.

\bibitem{macariorat2013}
G. Macario-Rat, J. Pl\^ut, and H. Gilbert,
\emph{New Insight into the Isomorphism of Polynomial Problem IP1S and Its Use in Cryptography},
IACR Cryptology ePrint Archive, Report 2013/799 (2013).
\href{https://eprint.iacr.org/2013/799}{ePrint 2013/799}.

\end{thebibliography}

\section*{AI-assisted research and writing disclosure}

The author used OpenAI Codex for exploratory computation, proof checking, and
language editing. The author reviewed the mathematical claims and assumes
responsibility for the manuscript.

\end{document}
